% Suggested text for the XDS Letter.
% The main-text section is intentionally short; methodological details are
% deferred to the Appendix.

\subsection{A detectability sweet spot}
\label{sec:sweetspot}

The detectability of the XDS is not controlled by the satellite alone, but by
the combined atmospheric, rotational, and dynamical properties of the
planet--moon system.  To identify the most favourable regime, we performed an
exploratory search over synthetic systems satisfying simple Roche, Hill-stability,
and rotational-breakup constraints (Appendix~\ref{app:sweetspot};
Fig.~\ref{fig:xds-sweetspot}).  The search
shows that high-S/N XDS signals preferentially arise for inflated, low-density
giant planets around relatively small stars, with extended H/He atmospheres,
moderate projected rotation, and close satellites.  In the adopted observing
setup, detectable configurations cluster around
$M_{\rm p}\simeq0.1$--$0.4\,M_{\rm J}$, $R_{\rm p}\gtrsim1.2\,R_{\rm J}$,
$\rho_{\rm p}\lesssim0.15\,{\rm g\,cm^{-3}}$, $v_{\rm eq}\simeq6$--$15\,
{\rm km\,s^{-1}}$, and $a_{\rm m}\simeq2$--$4\,R_{\rm p}$.  The strongest
synthetic cases correspond to super-puffy planets, for which the large scale
height boosts the atmospheric transmission signal while the moon remains close
enough to imprint a coherent time--velocity distortion during transit.

This result should be interpreted as a detectability map rather than a
population prediction.  It suggests that cold long-period giants are not
automatically optimal for XDS searches despite their favourable Hill spheres;
instead, the most promising targets are low-gravity, inflated giants with long
enough transits and sufficiently rapid atmospheric rotation.

\begin{figure*}
\centering
\includegraphics[width=\textwidth]{xds_paper_snr_scatter_maps_lhs100k_ref1500_seed11.png}
\caption{
Exploratory detectability search for the Exomoon Doppler Shadow.
Grey points show dynamically allowed synthetic systems refined with the
forward model, while coloured points indicate configurations with
${\rm S/N}_{\rm XDS}\ge5$ under the adopted observing assumptions
($R=10^5$, continuum S/N per resolution element of 300, 200 spectral lines,
one transit, and a 10 ppm systematic floor).  Cyan circles highlight the
highest-S/N detectable systems.  Detectable configurations cluster around
inflated low-density planets, large atmospheric scale heights, moderate
rotational velocities, and close satellites.
}
\label{fig:xds-sweetspot}
\end{figure*}


\appendix
\section{Exploratory sweet-spot search}
\label{app:sweetspot}

We complemented the forward model with an exploratory detectability search
designed to identify where the XDS signal is most favourable.  The search is
not a population synthesis and does not assume an occurrence rate for moons.
Its purpose is to map the regions of parameter space in which the signal would
be detectable under explicit, reproducible assumptions.

We sampled synthetic systems in the parameters
$(M_{\rm p},R_{\rm p},R_\star,P_{\rm p},P_{\rm rot},a_{\rm m},R_{\rm m})$.
The sampled ranges were
$M_{\rm p}=0.08$--$2\,M_{\rm J}$,
$R_{\rm p}=0.75$--$1.75\,R_{\rm J}$,
$R_\star=0.45$--$1.05\,R_\odot$,
$P_{\rm p}=3$--$300\,{\rm d}$,
$P_{\rm rot}=3$--$24\,{\rm h}$,
$a_{\rm m}=1.8$--$12\,R_{\rm p}$, and
$R_{\rm m}=0.12$--$0.35\,R_{\rm p}$.  Logarithmic sampling was used for
$M_{\rm p}$, $P_{\rm p}$, and $P_{\rm rot}$; the remaining quantities were
sampled linearly.  To improve coverage of the multidimensional space we used a
Latin-hypercube design rather than independent random draws.

Each system was first passed through a fast physical screen.  We required the
moon to lie outside the larger of the fluid Roche limit and the hard geometric
limit $1+R_{\rm m}/R_{\rm p}$, and inside $0.35\,R_{\rm H}$, where
$R_{\rm H}$ is the planetary Hill radius.  We also required the planetary spin
period to exceed $1.15$ times the breakup period.  The transit duration was
estimated for a central circular transit and systems with
$T_{\rm dur}<1.5\,{\rm h}$ were discarded.  The equilibrium temperature was
estimated assuming zero albedo, full redistribution, $M_\star=0.8\,M_\odot$,
and $T_\star=5200\,{\rm K}$; systems with $T_{\rm eq}<250\,{\rm K}$ were
excluded from this detectability exercise.

For the atmospheric scaling, we used
\begin{equation}
H=\frac{k_{\rm B}T}{\mu m_{\rm H}g},
\end{equation}
with $\mu=2.3$, and approximated the effective line depth as
\begin{equation}
\delta_{\rm atm}\simeq
\frac{2R_{\rm p}N_HH}{R_\star^2},
\end{equation}
where $N_H=5$ effective scale heights.  Candidate systems were ranked using a
fast score proportional to the atmospheric line depth, the projected lunar
area, a rotational Doppler factor, and a broad timing window in
$P_{\rm m}/T_{\rm dur}$.  This score was used only to select systems for
refinement; final S/N values were computed with the full time--velocity forward
model.

In the fiducial run we screened $10^5$ Latin-hypercube samples.  Of these,
$26954$ satisfied the fast physical filters.  We then refined the 1500
highest-scoring systems with the full XDS forward model, adopting
$R=10^5$, continuum S/N per resolution element of 300, $N_{\rm lines}=200$,
one transit, and a systematic floor of 10 ppm.  The matched-filter S/N was
computed as
\begin{equation}
{\rm S/N}_{\rm XDS} =
\left[\sum_{t,v}\left(\frac{\Delta F_{\rm XDS}(t,v)}{\sigma}\right)^2
\right]^{1/2}
\sqrt{N_{\rm lines}N_{\rm tr}},
\end{equation}
including the finite-exposure smearing penalty described in
Appendix~\ref{app:observability}.

The refined sample contained 268 systems with ${\rm S/N}\ge5$, 72 with
${\rm S/N}\ge10$, and 11 with ${\rm S/N}\ge20$.  The highest-S/N synthetic
case reached ${\rm S/N}=83.2$ and had
$M_{\rm p}=0.138\,M_{\rm J}$,
$R_{\rm p}=1.75\,R_{\rm J}$,
$\rho_{\rm p}=0.034\,{\rm g\,cm^{-3}}$,
$R_\star=0.49\,R_\odot$,
$T_{\rm eq}=434\,{\rm K}$,
$v_{\rm eq}=10.1\,{\rm km\,s^{-1}}$,
$a_{\rm m}=2.17\,R_{\rm p}$, and
$P_{\rm m}/T_{\rm dur}=20.6$.  These values are not proposed as a typical
system.  Rather, they identify the physical direction in which the XDS becomes
most detectable: very low surface gravity, a large atmospheric scale height,
moderate rotational broadening, and a close satellite whose shadow remains
coherent over the transit.

The main limitation of this exercise is astrophysical plausibility.  The most
favourable systems are super-puffy planets, and future work should impose
stronger constraints from atmospheric escape, tidal evolution, moon survival,
and realistic molecular line lists before converting this detectability map
into a target ranking.
